\section{Packet-level evaluation in ns-3}

The analytical model used during training assumes that offered demand is carried, which is
precisely the assumption that fails once a link saturates. Every reported result is therefore
produced by a discrete-event packet simulation in \mbox{ns-3} (v3.42), which acts as the judge of
the comparison rather than the analytical surrogate. The trained policy is evaluated in
\emph{export--simulate} form: for each held-out matrix the deterministic rollout
$a = \arg\max_a \pi_\theta(a\mid o,\mathbf{m})$ emits an explicit node path $\pi_i$ per demand,
and OSPF emits its shortest path $\pi_i^{\mathrm{sp}}$ for the same demand. Both routings are
written to a single JSON file and installed in \mbox{ns-3} as \emph{static} per-hop routes,
bypassing the global routing helper, so that the simulator reproduces the intended paths exactly
and the only difference between the two runs is the routing decision itself.

The scenario is built from the same topology JSON used everywhere else. Each link is a
point-to-point device with data rate $c_e$ and propagation delay taken from the instance; each
demand $i$ with $r_i > 0$ becomes a one-way constant-rate UDP \texttt{OnOff} application with
$1400$\,B payload sending at $r_i$, so that congestion manifests as genuine queue overflow and
loss rather than as transport back-off. To keep large topologies tractable, both rates and
capacities are divided by a common factor $\kappa$,
\begin{equation}
c_e \mapsto c_e/\kappa, \qquad r_i \mapsto r_i/\kappa
\quad\Longrightarrow\quad
u_e = 100\,\frac{\ell_e}{c_e} \;\; \text{invariant},
\end{equation}
which preserves every utilisation ratio, and hence loss and queueing behaviour, while reducing
the packet count by $\kappa$; $\kappa = 20$ and a simulated horizon of $8$\,s are used throughout.
On the largest instance the demand set is additionally restricted to the $N_f$ largest flows,
applied identically to OSPF and to the learned policy so that both are always compared on the
same flow set.

Three quantities are recovered from the simulation. Loss, mean end-to-end delay and delivered
throughput are collected with \texttt{FlowMonitor} over all flows,
\begin{equation}
\mathrm{loss} = 100\,\frac{P_{\mathrm{tx}} - P_{\mathrm{rx}}}{P_{\mathrm{tx}}}\,[\%],
\qquad
\bar{d} = \frac{1}{P_{\mathrm{rx}}}\sum_{p} \bigl(t^{\mathrm{rx}}_p - t^{\mathrm{tx}}_p\bigr),
\qquad
\Theta = \frac{8\,B_{\mathrm{rx}}}{W},
\end{equation}
with $P_{\mathrm{tx}}, P_{\mathrm{rx}}$ the transmitted and received packet counts,
$B_{\mathrm{rx}}$ the delivered bytes and $W$ the measurement window. Link utilisation is measured
from the bytes actually enqueued on each device rather than from the analytical load, giving the
packet-level bottleneck
\begin{equation}
u^{\mathrm{ns3}}_e = \frac{100}{\chi}\cdot\frac{8\,B^{\mathrm{enq}}_e}{c_e\,W},
\qquad
U^{\mathrm{ns3}}_{\max} = \max_{e\in\mathcal{E}} u^{\mathrm{ns3}}_e ,
\end{equation}
where $\chi = 7/8$ corrects for the fact that the applications are active for $7$\,s of the $8$\,s
accounting window. Note that $u^{\mathrm{ns3}}_e$ saturates at $100\%$ by construction: offered
load beyond line rate does not appear as utilisation above capacity but as queue drops, so an
analytical bottleneck of, say, $140\%$ surfaces in \mbox{ns-3} as a fully utilised link plus a
loss figure. Utilisation and loss must therefore be read together.

Results are stratified by congestion regime rather than pooled. A held-out matrix is labelled
\emph{overload} if OSPF cannot carry it feasibly and \emph{feasible} otherwise,
\begin{equation}
\mathrm{regime}(\mathbf{r}) =
\begin{cases}
\text{overload}, & U^{\mathrm{OSPF}}_{\max}(\mathbf{r}) > 100\%,\\[2pt]
\text{feasible}, & \text{otherwise},
\end{cases}
\end{equation}
and matrices are sampled to populate both regimes, since the two ask different questions: in the
feasible regime a learned policy can at best match OSPF on loss and should not inflate delay,
whereas the overload regime is where re-routing away from the bottleneck can convert drops into
delivered traffic. For each metric $m$ and regime, results are averaged over the held-out
matrices and independent training seeds and reported as $\mu \pm \sigma$,
\begin{equation}
\mu_m = \frac{1}{|\mathcal{S}||\mathcal{H}|}\sum_{s\in\mathcal{S}}\sum_{\mathbf{r}\in\mathcal{H}}
m\bigl(\pi_{\theta_s}, \mathbf{r}\bigr),
\end{equation}
so that seed-to-seed variance of the learned policy is visible in every reported figure rather
than hidden behind a single best run. OSPF, the myopic greedy best-response and the learned
policies are all simulated on identical matrices, identical flow sets and identical scenario
parameters, differing only in the installed paths.
